% -------------------------
\chapter{Literature Review}




\vspace{1cm}
\section{Existing Studies and Research Papers}

The modelling of time-varying volatility in financial markets has been a central concern of financial econometrics since the introduction of the Autoregressive Conditional Heteroskedasticity (ARCH) introduced by the \cite{engle1982}. Engle demonstrated that the variance of financial returns is not constant over time but follows a predictable pattern conditional on past information — a phenomenon now widely referred to as volatility clustering. The ARCH model differs from the earlier models that assumed the constant variance, providing a formal framework for capturing the empirical observation that large changes in asset prices tend to be followed by further large changes, and small changes by small changes (Mandelbrot, 1963).\\
specifically was later introduced by the \cite{bollerslev1986} addressing the limitation of ARCH model which required the large number of lagged terms to capture the persistence of volatility.  In the GARCH(1,1) current conditional variance depends on one lagged squared residual and one lagged conditional variance. GARCH(1,1) has since become the benchmark model in financial volatility research . Its widespread adoption reflects both its empirical performance and its theoretical elegance, as noted by \cite{hansen2005}  who compared 330 ARCH-type models and concluded that no ARCH model in his study consistently outperformed the simple GARCH(1,1) in forecasting exchange rate volatility.\\
But the GARCH model also have its limitation. GARCH model give symmetric response to positive and negative return shocks of equal magnitude.  This assumption has been consistently rejected in equity market data. Black (1976) first documented the leverage effect, whereby a decline in stock prices increases financial leverage and consequently raises future volatility by more than an equivalent price increase. To explain this asymmetry, Nelson (1991) introduced the Exponential GARCH (EGARCH) model, which allows negative shocks to have a disproportionately larger impact on conditional variance than positive shocks of the same size. Subsequently, \cite{glosten1993} introduced the GJR-GARCH model — also known as the Threshold GARCH — which captures the leverage effect through an additional indicator variable that activates only when return innovations are negative. Both models have been widely applied in equity market research and have consistently outperformed the symmetric GARCH specification in markets characterised by strong leverage effects \cite{ding1993}.\\
A further limitation of standard GARCH models is their assumption that volatility shocks decay at an exponential rate, implying that the effect of any given shock on future volatility is relatively short-lived. Empirical evidence, however, suggests that volatility in financial markets exhibits long memory — that is, the impact of shocks persists over much longer horizons than exponential decay would predict.\cite{baillie1996} addressed this by developing the Fractionally Integrated GARCH (FIGARCH) model, which introduces a fractional differencing parameter d to capture the slow hyperbolic decay of volatility persistence. When d equals zero, the FIGARCH model reduces to the standard GARCH; when d equals one, it becomes the Integrated GARCH (IGARCH). Values of d between zero and one indicate the presence of long memory, making FIGARCH particularly appropriate for markets where volatility shocks are slow to dissipate.



\section{Empirical Evidence from developed markets}

A lot of the academic literature has provided evidence for developed markets that, at least for some markets, some models of the GARCH-family tend to work better than others over different market conditions. For instance, Bollerslev (1992) was one of the earliest to consider ARCH-type models from an angle and it demonstrated that these could apply to a variety of market domains, including bond, foreign exchange and stock markets in both the United States and Europe. It additionally exposed the existence of volatility persistence in financial return series.It also demonstrated the volatility persistence of financial return series.

Hansen conducted an important study in 2005 which analyzed 330 models of the ARCH type, based on stock returns data from the Deutsche Mark and the IBM. The study found that no model was consistently better than the GARCH model when it came to forecasting exchange rate volatility although some models that took into account asymmetry did a bit better with stock data. In fact, when forecasting, it is not necessarily true that a complex model can give better forecasts; this is one of the reasons why we are comparing models in this study.

However, on its own history a study by Ding in 1993 revealed long-range dependence in the absolute and squared returns to the S\&P 500 index over 17 years. The FIGARCH model was developed by Baillie in 1996 and this study gave an indication for applying this model. Other studies, such as the one by Andersen in 1997 found that standard GARCH models do not capture volatility persistence in stock markets well which is another reason to use fractionally integrated models.

A big survey of 93 studies on volatility forecasting was done by Poon in 2003. It found that models based on historical volatility (models in the GARCH family) tended to have a higher forecasting performance than models based on implied volatility in terms of short-term forecasting. But it was also observed in the study that these different GARCH models perform better at different markets and time periods. This is why we have so many models and periods being explored in this study; no one model or time period is the best in all scenarios.


\section{Volatility Modelling in Indian Equity Markets}
The application of GARCH-family models to Indian equity markets has gained considerable momentum over the past two decades, driven by the rapid growth of the National Stock Exchange (NSE) and the Bombay Stock Exchange (BSE) as significant participants in global capital markets. However, despite this growing interest, the empirical evidence on the relative performance of competing volatility models in the Indian context remains mixed and inconclusive, providing the primary motivation for the present study.\\
Early applications of GARCH models to Indian equity markets focused predominantly on the benchmark Sensex and Nifty 50 indices using relatively short sample periods. \cite{kaur2004} was among the first to document the presence of volatility clustering and ARCH effects in BSE Sensex returns, confirming that the statistical preconditions for GARCH modelling are satisfied in the Indian market context. The study found significant persistence in conditional volatility, with the sum of ARCH and GARCH coefficients approaching unity — a finding replicated consistently in subsequent Indian market studies and corroborated by the present study's persistence estimates.\\
The leverage effect — the asymmetric response of volatility to positive and negative return shocks — has been documented extensively in Indian equity markets. \cite{pandey2005} applied GARCH, EGARCH, and GJR-GARCH models to Nifty 50 returns and found significant evidence of the leverage effect, with negative shocks generating disproportionately larger increases in conditional volatility than equivalent positive shocks. This finding was subsequently confirmed by \cite{banerjee2006}, who reported that asymmetric GARCH specifications consistently outperformed the symmetric GARCH model on BSE data in terms of in-sample fit, a result consistent with the AIC and BIC rankings reported in the present study.\\
The question of which specific GARCH model best captures Indian market volatility has attracted considerable attention but has not yielded a consensus. \cite{goudarzi2011} compared GARCH and EGARCH on BSE 500 data and found that EGARCH provided superior in-sample fit — consistent with the present study's information criteria results which rank EGARCH first among the four models estimate..\\
 \cite{mahajan2022} applied GARCH and RNN models to Nifty 50 data over the period 2014–2021 and found that GARCH-family models effectively captured the extreme volatility observed during the COVID-19 pandemic of 2020, though they noted that the post-crisis period exhibited distinctly different volatility dynamics from the pre-crisis period. This observation directly motivates the sub-period analysis conducted in the present study, which formally tests whether the COVID-19 pandemic constituted a structural break in Nifty 50 volatility dynamics.\\
A notable limitation of the existing Indian market literature is its tendency to focus on either in-sample model fit or out-of-sample forecasting in isolation, rarely evaluating both dimensions simultaneously on the same dataset. Furthermore, the majority of existing studies employ sample periods that either predate the COVID-19 pandemic or cover only the immediate crisis period, leaving the post-recovery dynamics of Indian market volatility largely unexplored. The present study addresses both limitations by combining a rigorous in-sample model comparison with a comprehensive multi-horizon, multi-quarter out-of-sample forecasting evaluation over the full year 2025 — a period that has received no attention in the published literature to date.


\section{Pre and post-COVID volatility studies}

The COVID-19 pandemic was a surprise to global financial markets. From January to March 2020 the Nifty 50 index went down by 38 percent, which is one of the biggest drops in Indian equity markets history and even bigger than what happened during the 2008 Global Financial Crisis in some countries.

Some researchers, Baker and others found that the US equity markets were very volatile during this time. They said this was because people were not sure what would happen to the economy when everything was locked down. Another study by Zhang and others looked at ten stock markets and found that they all had a big increase in volatility especially in countries like India.

There is a gap in what we know about the Indian market. Most studies only look at how volatility models work inside the sample or how well they predict outside the sample but not both. Also most studies only look at what happened before the pandemic or during the pandemic and not what happened after. This study tries to fill this gap by looking at four GARCH models. How well they work inside the sample from 2010 to 2024 and then seeing how well they predict what will happen in the four quarters of 2025. The COVID-19 pandemic and its effects on the market are what we are focusing on. The Nifty 50 index and its big drop are important to understand. The COVID-19 pandemic had an impact, on global financial markets, including the Indian market.

The effects of the epidemics, on the 50 and market volatility, are investigated here.The pandemics impact, on the 50 and market volatility, are studied here. persistence than developed markets during the crisis period. These findings established the theoretical expectation that COVID-19 would leave a lasting imprint on the volatility structure of emerging market indices such as the Nifty 50.\\
In the Indian context, several studies specifically examined the impact of COVID-19 on Nifty 50 volatility using GARCH-family models. \cite{mishra2020} applied GARCH(1,1) and EGARCH models to daily Nifty 50 returns and found a significant increase in volatility persistence during the pandemic period, with the ARCH and GARCH coefficients rising substantially compared to the pre-pandemic estimation window. This finding is consistent with the theoretical expectation that external shocks increase the sensitivity of conditional variance to new information, reflected in higher alpha coefficients in the post-shock period.\\
The post-COVID recovery period has received comparatively less attention in the Indian market literature. \cite{rai2022} examined Nifty 50 volatility across pre-COVID, during-COVID, and post-COVID sub-periods and found that while volatility gradually declined from its 2020 peak, it did not fully revert to pre-pandemic levels by the end of their sample period in mid-2021. This finding of partial mean reversion is particularly relevant to the present study, which extends the analysis through December 2024 and therefore captures the longer-term trajectory of post-COVID volatility normalisation.\\
A common limitation across the existing COVID-era studies is their focus on narrow sample windows that capture only the immediate crisis period or the early recovery phase. Furthermore, most studies estimate a single GARCH specification across the entire sample without formally testing for structural breaks, which risks parameter contamination . The present study addresses both limitations by employing a Bai-Perron structural break test to formally identify data-driven regime boundaries, and by estimating all four GARCH models separately on three distinct regimes — the pre-COVID period (January 2010 to February 2020), the COVID crisis and recovery period (February 2020 to May 2022), and the post-crisis period (May 2022 to December 2024). This design allows for a more nuanced assessment of how the pandemic permanently altered the volatility dynamics of the Indian equity market, contributing novel empirical evidence to a literature that has largely been confined to shorter and less methodologically rigorous comparative frameworks.

\section{Summary and Research Gap}

The preceding sections have reviewed the theoretical and empirical literature on volatility modelling across four dimensions — the foundational development of GARCH-family models, their empirical validation in developed markets, their application to Indian equity markets, and the emerging literature on COVID-19 induced volatility regime shifts. This section synthesises the key findings from this review and identifies the specific gaps that the present study seeks to address.\\
The review of the theoretical literature in Section 2.1 established that GARCH-family models represent the dominant framework for modelling conditional volatility in financial markets, with the symmetric GARCH(1,1) of \cite{bollerslev1986} serving as the benchmark specification against which more complex models are evaluated. The asymmetric extensions — EGARCH \cite{nelson1991conditional} and \cite{glosten1993}GJR-GARCH  were developed to capture the well-documented leverage effect, while FIGARCH \cite{baillie1996} addresses the long-memory properties of volatility that standard specifications fail to accommodate. Each model therefore captures a distinct dimension of volatility dynamics, and the selection of the most appropriate specification for a given market and time period is fundamentally an empirical question.\\

The review of evidence in Sections 2.2 and 2.3 shows that GARCH-family models have been used a lot in both Indian equity markets.. The existing literature has three important limitations. First most Indian market studies only use one or two GARCH specifications, like GARCH(1,1). Egarch. They do not compare all the types of symmetric, asymmetric and long-memory specifications. This makes it hard to say which models are better for the market.

Second existing studies focus much on how well the models fit the data but they do not check how well the models forecast what will happen in the future. The study by Hansen in 2005 showed that just because a model fits the data well it does not mean it will forecast well.. This has not been looked at much in the Indian GARCH literature.

Third most studies use GARCH models on a sample of data without checking if there are any big changes in the data. This can cause problems when the sample covers different periods of volatility.

The COVID-19 literature reviewed in Section 2.4 shows a gap. There are not studies that look at how the Indian market volatility changed from before the pandemic to during the pandemic and after the pandemic. Some studies show that volatility was higher during the pandemic. They do not say if it went back to normal after the pandemic. Also the literature does not look much at how the Indian equity market volatility changed after the pandemic with the geopolitical and monetary policy shocks in 2022.

This study addresses all four of these gaps. First it compares four GARCH-family specifications. GARCH, EGARCH, GJR-GARCH and FIGARCH. Using daily Nifty 50 returns from January 2010 to December 2024. This is the comprehensive comparison of GARCH models on Indian equity market data so far.

Second it checks how well the models forecast what will happen in the future over four periods. 1-Day 5-day 10-day and 20-day. It uses different metrics to measure accuracy, including MAE, RMSE and the Diebold-Mariano test. This gives an idea of how well the models work outside the sample of data.

Third it uses the Bai-Perron structural break test to find the boundaries between volatility regimes. This is better than dividing the data into arbitrary periods.

Fourth it looks at volatility dynamics over three periods. Before COVID during the COVID crisis and recovery and after COVID. This gives a picture of how the pandemic and other shocks changed the Indian market volatility.

This study makes three contributions to the literature. It provides the comprehensive study of volatility forecasting on Nifty 50 data so far. It gives evidence on how COVID-19 changed the Indian equity market volatility dynamics.. It shows that the ranking of models based on how well they fit the data is different from the ranking based on how well they forecast what will happen in the future, which is important for risk management, in the Indian market.

